\section{Verwendete Modelle}

Die Architektur der Datenpipeline erlaubt es, verschiedene Modelle zu testen.
Es werden sowohl Baseline-Modelle als auch komplexere Modelle trainiert.
Die Baseline-Modelle dienen als Referenz, um die Leistung der komplexeren Ansätze besser einordnen zu können.

\subsubsection{Naive Baseline}

Die naive Baseline liefert als Vorhersage den Wert der letzten Eingabe zurück \cite{Baseline2024}.

\subsubsection{Mean Baseline}

Die Mean Baseline liefert, wie der Name bereits andeutet, den Mittelwert der Trainingsdaten als Vorhersage \cite{Baseline2024}.

\subsubsection{Lineare Regression}

Die lineare Regression ist ein Standardverfahren, um lineare Zusammenhänge in Daten vorherzusagen.
Die Daten werden durch eine Gerade angenähert, auf deren Basis neue Werte prognostiziert werden können \cite{Linear2026}.

\subsubsection{Ridge Regression}

Die Ridge Regression ist eine Erweiterung der linearen Regression um L2-Regularisierung, die die Größe der Koeffizienten minimiert und so Overfitting verhindern soll \cite{Ridge2026}.


\subsubsection{XGB Regressor}

XGB (Extreme Gradient Boosting) ist ein Regressionsmodell, das Vorhersagen mit einer Folge von Entscheidungsbäumen trifft.
Durch die Kombination mehrerer Bäume können sowohl nichtlineare Zusammenhänge als auch komplexe Zeitreihenstrukturen modelliert werden \cite{XGBoost}.


\subsubsection{LGBM Regressor}

LGBM (Light Gradient Boosting Machine) ist eine auf Geschwindigkeit optimierte Version von XGB.
Es ist bei großen Datenmengen effizienter als XGB \cite{LightgbmLGBMRegressor}.

\subsubsection{SVR}

SVR (Support Vector Regressor) ist eine Variante der SVM (Support Vector Machine) für Regressionsaufgaben.
Das Verfahren ist robust gegenüber Overfitting und kann komplexe Muster erfassen, weshalb es häufig für die Prognose von Werten verwendet wird \cite{SVR}.

\subsubsection{KNN}

KNN (K-Nearest Neighbors Regressor) liefert Vorhersagen auf Grundlage ähnlicher Punkte in den Trainingsdaten.
Dazu müssen zunächst die Vergleichswerte (Lag-Werte) aus den Trainingsdaten erzeugt werden.
Für Vorhersagen werden diese Werte dann mit den ähnlichsten Beobachtungen verglichen und daraus der Mittelwert gebildet \cite{Knearest2026}.

\subsubsection{LSTM}

Ein LSTM-Netzwerk (Long Short-Term Memory) ist ein von Sepp Hochreiter und Jürgen Schmidhuber im Jahr 1997 entwickeltes Verfahren \cite{Long2026}, das im Gegensatz zu anderen rekurrenten neuronalen Netzwerken ein Gedächtnis über längere Zeiträume modellieren kann.
Dies geschieht durch drei Arten von Toren: das Eingangstor, das Vergessens- bzw. Merktor sowie das Ausgangstor.
Diese ermöglichen bei der Rückführung des Fehlersignals eine effektivere Aktualisierung der Gewichte.

Aufgrund seiner Speicherfähigkeit eignet sich ein LSTM-Netzwerk besonders gut zum Lernen von Zeitreihen und wird daher in dieser Arbeit als zentrales Modell für die Vorhersage eingesetzt \cite{Long2026}.



